% GENERATED FILE. Edit canonical lessons, then run npm run generate:books.
% canonical-source-hash: fnv1a64:dfde41b24194b059
% canonical-lessons: ZH-C38-chi, ZH-W38-chi, ZH-C38-fan, ZH-W38-fan, ZH-C38-mifan, ZH-C38-e, ZH-C38-cai

\chapter{Eating}
\label{ch:zh-eating}
\hlchaptermodality{38}

\begin{chapteropening}
\textbf{By the end of this chapter:} \emph{I can say what I eat and ask what somebody likes to eat.}

Complete the last lesson of chapter 38: I can say what I eat and ask what somebody likes to eat.
\end{chapteropening}

\section[\texorpdfstring{\zh{吃} (chī)}{吃 (chī)}]{\zh{吃} (chī) — to eat}

\label{lesson:ZH-C38-chi}



\begin{quote}
Before the new one: say the Mandarin for which one, then the Mandarin for that one.
\end{quote}

\subsection*{You'll want to know: \zh{吃}}
\textbf{\zh{吃}} — \emph{chī} — \textquotedblleft{}to eat\textquotedblright{}.

\begin{quote}
\textbf{\zh{你吃什么}} \emph{nǐ chī shénme} — What are you having?
\end{quote}

The new character, in pieces: \zh{吃} — \zh{口}, the mouth piece, on the left, and on the right three strokes: a short fall, a short bar, and one long stroke that runs across, drops and hooks up. The next lesson puts it in your hand.

\begin{cousinweb}
\textbf{\zh{吃}} (\emph{chī}) is probably inherited from Middle Chinese \emph{khek}, which the Guangyun glosses 'to eat' and writes with a different character; there \zh{吃} itself is \emph{kj+t}, 'to stammer' (Guangyun, in Baxter's transcription). How the word and this character came together is not settled.
\end{cousinweb}

\begin{grammarlens}[title={What you've built}]
One more word for everyday answers.
\end{grammarlens}

\subsection*{Your turn}
\begin{itemize}
\raggedright
  \item \emph{Say it:} \emph{chī}
  \item \emph{Say it:} \emph{chī}, once more
  \item \emph{Say it:} \emph{nǐ chī shénme} — What are you having?
  \item \emph{Recall:} say the Mandarin for which one, then the Mandarin for that one, then say \emph{chī} again
\end{itemize}

\begin{tcolorbox}[breakable,colback=teal!4,colframe=teal!35!black,title={Before you move on}]
What does \zh{吃} mean? (\textquotedblleft{}To eat\textquotedblright{}.) Say it once more.
\end{tcolorbox}
\section[\texorpdfstring{\zh{吃} (chī)}{吃 (chī)}]{\zh{吃} — a mouth, then three strokes}

\label{lesson:ZH-W38-chi}



\begin{quote}
The word first, then the character inside it.

\begin{itemize}
\raggedright
  \item \emph{Recall:} say \emph{chī} — \textbf{R1}, one lesson back, before it decays
\end{itemize}

\begin{quote}
\zh{口} \zh{叫}
\end{quote}

\emph{Read it:} each one aloud, with nothing said first — \emph{kǒu}, \emph{jiào}
\end{quote}

\subsection*{Script — a mouth, then three strokes}
\hlblockfigure{\detokenize{figures/ZH-W38-chi-filmstrip.pdf}}{How \zh{吃} is written, stroke by stroke}

\begin{quote}
\zh{吃}
\end{quote}

\textbf{Two pieces, left before right.}

\textbf{The left piece, \zh{口}, three strokes}, small, as in \zh{叫}.

\textbf{The right piece, three strokes.} A short left-falling stroke. A short horizontal under it. Then one long stroke: across, down, bending right along the bottom and hooking up.

3 + 3 = six.

\subsection*{Your turn}
\begin{itemize}
\raggedright
  \item \emph{Write it:} \zh{吃}, all six, with the model in front of you
  \item \emph{Cover:} write it again from memory, and count the pieces, not the strokes
\end{itemize}

\begin{tcolorbox}[breakable,colback=teal!4,colframe=teal!35!black,title={Before you move on}]
How many pieces in \zh{吃}? (\textbf{Two} — \zh{口}, then three strokes.) Which character from the last chapters also has a small \zh{口} on the left? (\textbf{\zh{叫}}.) How does the last stroke end? (\textbf{It hooks up}.)

Source: \href{https://raw.githubusercontent.com/chanind/hanzi-writer-data/68d10a4b21150cae5e1ebbd223eed289cf32d90c/data/\%E5\%90\%83.json}{Hanzi Writer Data \zh{吃}, pinned stroke paths and medians}.
\end{tcolorbox}
\section[\texorpdfstring{\zh{饭} (fàn)}{饭 (fàn)}]{\zh{饭} (fàn) — cooked rice; meal}

\label{lesson:ZH-C38-fan}



\begin{quote}
Before the new one: say the Mandarin for that one, then the Mandarin for to eat.
\end{quote}

\subsection*{You'll want to know: \zh{饭}}
\textbf{\zh{饭}} — \emph{fàn} — \textquotedblleft{}cooked rice; meal\textquotedblright{}.

\begin{quote}
\textbf{\zh{你吃饭了吗}} \emph{nǐ chī fàn le ma} — Have you eaten?
\end{quote}

The new character, in pieces: \zh{饭} — a three-stroke food piece on the left, and a four-stroke piece on the right. The next lesson puts it in your hand.

\begin{cousinweb}
\textbf{\zh{饭}} (\emph{fàn}) is inherited from Middle Chinese \emph{bjonH} (Guangyun, in Baxter's transcription), glossed 'cooked grain'.
\end{cousinweb}

\begin{grammarlens}[title={What you've built}]
One more word for everyday answers.
\end{grammarlens}

\subsection*{Your turn}
\begin{itemize}
\raggedright
  \item \emph{Say it:} \emph{fàn}
  \item \emph{Say it:} \emph{fàn}, once more
  \item \emph{Say it:} \emph{nǐ chī fàn le ma} — Have you eaten?
  \item \emph{Recall:} say the Mandarin for that one, then the Mandarin for to eat, then say \emph{fàn} again
\end{itemize}

\begin{tcolorbox}[breakable,colback=teal!4,colframe=teal!35!black,title={Before you move on}]
What does \zh{饭} mean? (\textquotedblleft{}Cooked rice; meal\textquotedblright{}.) Say it once more.
\end{tcolorbox}
\section[\texorpdfstring{\zh{饭} (fàn)}{饭 (fàn)}]{\zh{饭} — the food piece, then four strokes}

\label{lesson:ZH-W38-fan}



\begin{quote}
The word first, then the character inside it.

\begin{itemize}
\raggedright
  \item \emph{Recall:} say \emph{fàn} — \textbf{R1}, one lesson back, before it decays
\end{itemize}

\begin{quote}
\zh{吃} \zh{了}
\end{quote}

\emph{Read it:} each one aloud, with nothing said first — \emph{chī}, \emph{le}
\end{quote}

\subsection*{Script — the food piece, then four strokes}
\hlblockfigure{\detokenize{figures/ZH-W38-fan-filmstrip.pdf}}{How \zh{饭} is written, stroke by stroke}

\begin{quote}
\zh{饭}
\end{quote}

\textbf{Two pieces, left before right.}

\textbf{The food piece, three strokes.} A left-falling stroke. A short horizontal that hooks down at its end. A vertical that flicks up to the right at the foot.

\textbf{The right piece, four strokes.} A short left-falling stroke along the top. A long left-falling stroke down the left. A horizontal that turns and falls to the left. Then a right-falling stroke across it.

3 + 4 = seven.

\subsection*{Your turn}
\begin{itemize}
\raggedright
  \item \emph{Write it:} \zh{饭}, all seven, with the model in front of you
  \item \emph{Cover:} write it again from memory, and count the pieces, not the strokes
\end{itemize}

\begin{tcolorbox}[breakable,colback=teal!4,colframe=teal!35!black,title={Before you move on}]
How many pieces in \zh{饭}? (\textbf{Two} — the food piece, then four strokes.) How does the food piece end? (\textbf{With a vertical that flicks up}.) Which stroke is last? (\textbf{The right-falling stroke}.)

Source: \href{https://raw.githubusercontent.com/chanind/hanzi-writer-data/68d10a4b21150cae5e1ebbd223eed289cf32d90c/data/\%E9\%A5\%AD.json}{Hanzi Writer Data \zh{饭}, pinned stroke paths and medians}.
\end{tcolorbox}
\section[mǐfàn]{mǐfàn — cooked rice}

\label{lesson:ZH-C38-mifan}



\begin{quote}
Before the new one: say the Mandarin for to eat, then the Mandarin for meal.
\end{quote}

\subsection*{You'll want to know: mǐfàn}
\emph{mǐfàn} — \textquotedblleft{}cooked rice\textquotedblright{}.

\begin{quote}
\emph{wǒ xǐhuan chī mǐfàn} — I like rice.
\end{quote}

Its characters come later; for now it is a word for the ear.

\begin{cousinweb}
\emph{mǐfàn} is built from mǐ 'husked grain' and \zh{饭} 'cooked grain', both inherited from Middle Chinese, \emph{mejX} and \emph{bjonH} (Guangyun, in Baxter's transcription).
\end{cousinweb}

\begin{grammarlens}[title={What you've built}]
One more word for everyday answers.
\end{grammarlens}

\subsection*{Your turn}
\begin{itemize}
\raggedright
  \item \emph{Say it:} \emph{mǐfàn}
  \item \emph{Say it:} \emph{mǐfàn}, once more
  \item \emph{Say it:} \emph{wǒ xǐhuan chī mǐfàn} — I like rice.
  \item \emph{Recall:} say the Mandarin for to eat, then the Mandarin for meal, then say \emph{mǐfàn} again
\end{itemize}

\begin{tcolorbox}[breakable,colback=teal!4,colframe=teal!35!black,title={Before you move on}]
What does mǐfàn mean? (\textquotedblleft{}Cooked rice\textquotedblright{}.) Say it once more.
\end{tcolorbox}
\section[è]{è — hungry}

\label{lesson:ZH-C38-e}



\begin{quote}
Before the new one: say the Mandarin for meal, then the Mandarin for cooked rice.
\end{quote}

\subsection*{You'll want to know: è}
\emph{è} — \textquotedblleft{}hungry\textquotedblright{}.

\begin{quote}
\emph{wǒ hěn è} — I'm starving.
\end{quote}

Its characters come later; for now it is a word for the ear.

\begin{cousinweb}
\emph{è} is inherited from Middle Chinese \emph{ngaH} (Guangyun, in Baxter's transcription), glossed 'not full'.
\end{cousinweb}

\begin{grammarlens}[title={What you've built}]
One more word for everyday answers.
\end{grammarlens}

\subsection*{Your turn}
\begin{itemize}
\raggedright
  \item \emph{Say it:} \emph{è}
  \item \emph{Say it:} \emph{è}, once more
  \item \emph{Say it:} \emph{wǒ hěn è} — I'm starving.
  \item \emph{Recall:} say the Mandarin for meal, then the Mandarin for cooked rice, then say \emph{è} again
\end{itemize}

\begin{tcolorbox}[breakable,colback=teal!4,colframe=teal!35!black,title={Before you move on}]
What does è mean? (\textquotedblleft{}Hungry\textquotedblright{}.) Say it once more.
\end{tcolorbox}
\section[cài]{cài — dish of food; vegetable}

\label{lesson:ZH-C38-cai}



\begin{quote}
Before the new one: say the Mandarin for cooked rice, then the Mandarin for hungry.
\end{quote}

\subsection*{You'll want to know: cài}
\emph{cài} — \textquotedblleft{}dish of food; vegetable\textquotedblright{}.

\begin{quote}
\emph{nǐ xǐhuan chī shénme cài} — What sort of food do you like?
\end{quote}

Its characters come later; for now it is a word for the ear.

\begin{cousinweb}
\emph{cài} is inherited from Middle Chinese \emph{tshojH} (Guangyun, in Baxter's transcription), which the dictionary explains as any plant that can be eaten.
\end{cousinweb}

\begin{grammarlens}[title={What you've built}]
One more word for everyday answers.
\end{grammarlens}

\subsection*{Your turn}
\begin{itemize}
\raggedright
  \item \emph{Say it:} \emph{cài}
  \item \emph{Say it:} \emph{cài}, once more
  \item \emph{Say it:} \emph{nǐ xǐhuan chī shénme cài} — What sort of food do you like?
  \item \emph{Recall:} say the Mandarin for cooked rice, then the Mandarin for hungry, then say \emph{cài} again
\end{itemize}

\begin{tcolorbox}[breakable,colback=teal!4,colframe=teal!35!black,title={Before you move on}]
What does cài mean? (\textquotedblleft{}Dish of food; vegetable\textquotedblright{}.) Say it once more.
\end{tcolorbox}
